\subsubsection{数据预处理}

问题二使用附件2的全年实测序列，负载与光伏各含365天、每天144个时段的功率观测。预处理围绕功率折算、时间对齐、净负荷合成和信息截断展开。

第一步是功率到能量的折算。与问题一一致，我们仍以10分钟为一个时段，把附件2的两列功率序列折算为时段能量：
\begin{equation}\label{eq:src12}
    L_{t}^{(d)}=P_{\text{load}}^{(d)}(t)\cdot \Delta t,\qquad V_{t}^{(d)}=P_{\text{pv}}^{(d)}(t)\cdot \Delta t
\end{equation}
我们将所有功率统一折算为时段电量，日序号与时段序号沿用符号说明中的定义。

第二步是净负荷的合成与信息截断。储能既可补偿缺口，也可吸收光伏余电，故我们保留净负荷的符号：
\begin{equation}\label{eq:src13}
    N_t^{(d)}=L_t^{(d)}-V_t^{(d)},\qquad \hat N_t^{(d)}=\hat L_t^{(d)}-\hat V_t^{(d)}.
\end{equation}
$N_t^{(d)}$用于实际执行与事后误差评价，$\hat N_t^{(d)}$仅由决策日前历史构造。$\alpha=0.8$表示日前保护的误差分位数水平，当天原计划提交后按约定电量结算。

第三步是同时段基准序列的构造。为在不使用当日实际值的前提下刻画日内形态，我们取目标日之前最近7个完整日的实测值构造加权同时段基准：
\begin{equation}\label{eq:src14}
    B_{t}^{(d)}=\sum_{k=1}^{K_d}\omega_{d,k}\,X_{t}^{(d-k)},\qquad
    \omega_{d,k}=\frac{K_d+1-k}{\sum_{j=1}^{K_d}j},\quad 1\le K_d\le7
\end{equation}
其中，$X$表示光伏序列，$K_d$为目标日前可用完整历史日数与7的较小值；无完整历史日时回退附件1的光伏基线。$k=1$为最近一天，权重从最近日的$K_d$逐日递减至最早日的1后归一化；满7天时即为$\omega_{d,k}=(8-k)/28$。负载则以最近一周同日同时段值为基准，再叠加历史周同比偏差的中位数。两类预测均不使用当日未来实际值。

第四步是预测偏差的滚动更新。令$e_t^{(d)}=N_t^{(d)}-\hat N_t^{(d)}$，按已完成日期保存同小时的六个10分钟误差。我们保留28日档案，但日前保护量只读取同小时最近28个误差观测：
\begin{equation}\label{eq:src15}
 B_{d,h}=\max\{0,Q_{0.8}(\mathcal R_{d,h})\},\qquad
 \widetilde N_t^{(d)}=\hat N_t^{(d)}+B_{d,h(t)}.
\end{equation}
至少积累8个完成日后我们启用保护；小时内样本不足时将保护量取零。问题三的日前保护改用最近5个完成日，日内备选方案的70\%、90\%分位数另行定义，不与本问的日前保护混用。

经上述处理，我们得到分位数保护后的日前输入，计划由滚动线性规划生成，执行阶段按实际缺口更新。


